%
% keplerreihe.tex -- Lösung der Kepler-Gleichung, E-M über M für drei
%                    Exzentrizitäten und die Bessel-Reihe für
%                    epsilon = 0.8, Daten aus kepler.dat (kepler.py)
%
% (c) 2026 Gian Cavegn, OST Ostschweizer Fachhochschule
%
\documentclass[tikz]{standalone}
\usepackage{amsmath}
\usepackage{times}
\usepackage{txfonts}
\usepackage{pgfplots}
\pgfplotsset{compat=1.18}
\usetikzlibrary{arrows,intersections,math,calc}
\definecolor{darkred}{rgb}{0.8,0,0}
\definecolor{darkgreen}{rgb}{0,0.6,0}
\begin{document}
\def\skala{1}
\begin{tikzpicture}[>=latex,scale=\skala]
\begin{axis}[
	width=11.4cm,
	height=8cm,
	xmin=0, xmax=6.2832,
	ymin=-0.9, ymax=0.9,
	xtick={0,1.5708,3.1416,4.7124,6.2832},
	xticklabels={$0$,$\pi/2$,$\pi$,$3\pi/2$,$2\pi$},
	ytick={-0.8,-0.4,0.4,0.8},
	axis lines=middle,
	axis line style={-},
	clip=false,
	grid=both,
	grid style={very thin,gray!25},
	legend pos=north east,
	legend cell align=left,
	tick label style={font=\small},
	xticklabel style={yshift=-2pt}
]
\addplot[thick,blue]      table[x=M,y=d02] {kepler.dat};
\addlegendentry{$\varepsilon=0.2$}
\addplot[thick,red]       table[x=M,y=d05] {kepler.dat};
\addlegendentry{$\varepsilon=0.5$}
\addplot[thick,darkgreen] table[x=M,y=d08] {kepler.dat};
\addlegendentry{$\varepsilon=0.8$}
% Bessel-Reihe fuer epsilon = 0.8, gestrichelt ein Glied, gepunktet drei Glieder
\addplot[thick,darkgreen,dashed]         table[x=M,y=s1] {kepler.dat};
\addlegendentry{$\varepsilon=0.8$ (1 Glied)}
\addplot[thick,darkgreen,densely dotted] table[x=M,y=s3] {kepler.dat};
\addlegendentry{$\varepsilon=0.8$ (3 Glieder)}
% Achsen über das Gitter hinaus verlängern, wie in den übrigen Abbildungen des Buches
\draw[->] (axis cs:6.2832,0) -- (axis cs:6.65,0) node[below] {$M$};
\draw[->] (axis cs:0,0.9) -- (axis cs:0,1.02) node[right] {$E-M$};
\end{axis}
\end{tikzpicture}
\end{document}
